% 自动生成，勿手改（gen_notes.py）
\chapter{译校批注}

本附录汇集正文中以源码注释形式给出的译校批注，按章节顺序排列。每条注明所在原文页码，便于对照原书。批注只写三类内容：易错预警、省略论证的点拨、背景直觉。\par\medskip

\section*{序与理论的概要与展望}

\noindent \textbf{（原书 p.xviii）}\ 本导论的叙述顺序（古典力学 → 复线丛与量子化 → sigma 模型 → Feynman 路径积分 → WZW/Chern--Simons）就是全书的概念骨架：第 1 章的共形块、第 2 章的辫群表示与 Witten 不变量、第 3 章的微扰展开都能挂在这条线上。初读可先记住 p.xv 的「态空间 = 复线丛的截面空间」与 p.xvii 的「无限维对象 → 有限维不变量」两句话。\par\smallskip

\noindent \textbf{（原书 p.xviii）}\ p.x 中 Dirac 的对应写的是 $\hat f_1\hat f_2-\hat f_2\hat f_1=-\sqrt{-1}\hbar\hat f_3$，而 p.x 末把 $\hat f$ 归一化为 $\sqrt{-1}\hbar\tilde f/\hbar$；两处相差一个常数因子，是作者为把 $\hbar$ 显式提出而作的约定，对易子 $[\hat f_1,\hat f_2]=-\sqrt{-1}\hbar\widehat{\{f_1,f_2\}}$ 与此一致。\par\smallskip

\noindent \textbf{（原书 p.xviii）}\ p.xi 的曲率恒等式 $\nabla_X\nabla_Y-\nabla_Y\nabla_X=\nabla_{[X,Y]}-2\pi\sqrt{-1}\omega(X,Y)$ 中，$2\pi\sqrt{-1}$ 的因子来自联络的约定 $\nabla=d-2\pi\sqrt{-1}\alpha$（第一 Chern 形式为 $d\alpha$ 而非 $2\pi$ 倍）。不同教材常把 $2\pi$ 并入 $\omega$，对照时须先确认此约定。\par\smallskip

\noindent \textbf{（原书 p.xviii）}\ p.xv 的三条公理只写到「粘合规则」，作者在此省略了「$\dim Z_Y<\infty$」与「$Z(X)$ 对 $X$ 的微分同胚不变」两条；Atiyah 的拓扑场论公理共五条，第 2 章 §2.4 会再用到。\par\smallskip

\noindent \textbf{（原书 p.xviii）}\ p.xvi 的 $CS(A)=\frac{1}{8\pi^{2}}\int_M\Tr(A\wedge dA+\frac{2}{3}A\wedge A\wedge A)$ 中，$\frac{2}{3}$ 的系数与 $\frac{1}{8\pi^{2}}$ 的归一化都和 $SU(2)$ 的 Killing 型取法有关；第 2 章 §2.5 与第 3 章 §3.2 沿用同一约定，跨章引用时不要换成别的归一化。\par\smallskip

\section*{共形场理论的几何}

\noindent \textbf{（原书 p.2）}\ 本书的记号前缀不可混：$LG$ 是环路群（映射空间），$L\mathfrak{g}=\mathfrak{g}\otimes\mathbb{C}((t))$ 是与之对应的环路代数（形式 Laurent 版本），$\widehat{\mathfrak{g}}=L\mathfrak{g}\oplus\mathbb{C}c$ 是它的中心扩张。后文 $LG_{\mathbb{C}}$ 指 $G_{\mathbb{C}}$ 的环路群，与 $\widehat{\mathfrak{g}}$ 不是同一层次的对象。\par\smallskip

\noindent \textbf{（原书 p.5）}\ 末句用 $[H\otimes t,H\otimes t^{-1}]=0$ 而余循环值 $\langle H,H\rangle\ne 0$，说明该扩张不能靠 $\mu([x,y])$ 消去，即上同调类非平凡；关键是括号为零而中心项不为零。\par\smallskip

\noindent \textbf{（原书 p.5）}\ 证明中的 $t$ 是 $\exp(tZ)$ 的参数，与形式变量 $t$ 无关，只是符号被重复使用；读 $gXg^{-1}=X+t[Z,X]+O(t^{2})$ 时不要把它当作 Laurent 变量。\par\smallskip

\noindent \textbf{（原书 p.5）}\ 命题 1.1 中 $df\,g$ 的 $d$ 只作用于 $f$；若写成 $d(fg)$ 或 $f\,dg$，因全微分的留数为零会差一个符号，$\mathrm{Res}(df\,g)=-\mathrm{Res}(f\,dg)$。\par\smallskip

\noindent \textbf{（原书 p.7）}\ 式 (1.6) 中中心元 $c$ 的系数是 $(m^{3}-m)/12$，即把中心荷归一到 1；文献里常见写法是 $(c/12)(m^{3}-m)$ 或 $(c/12)(m^{3}-m)\delta_{m+n,0}$，与此处「$c$ 本身是基元」的约定差一个倍数，换算时勿混。\par\smallskip

\noindent \textbf{（原书 p.8）}\ 式中的 $2\pi\sqrt{-1}$ 不能省成 $\sqrt{-1}$：正是这个因子使 $d\alpha$ 的 de Rham 类落在 $H^{2}(M,\Z)$ 的像里，即整性条件与 $2\pi\sqrt{-1}$ 的归一化绑定。\par\smallskip

\noindent \textbf{（原书 p.9）}\ $\pi_{1}(\Omega G)\cong\pi_{2}(G)$ 来自道路空间纤维化 $PG\to G$ 的同伦长正合列；这是「$LG$ 单连通」的全部依据，书中未展开。\par\smallskip

\noindent \textbf{（原书 p.10）}\ 式 (1.8) 的 $1/(24\pi^{2})$ 是使 $\sigma$ 在 $SU(2)\cong S^{3}$ 上积分为 1 的归一化，从而 $\sigma$ 的 de Rham 类是 $H^{3}(SU(2),\Z)$ 的整生成元；常数若换成 $1/8\pi^{2}$ 等便失去整性。\par\smallskip

\noindent \textbf{（原书 p.10）}\ $LG$ 单连通只保证 $H^{1}(LG,\Z)=0$，$H^{2}$ 一般不为零。事实上 $LG\cong G\times\Omega G$，故 $H^{2}(LG,\Z)\cong H^{2}(G,\Z)\oplus H^{2}(\Omega G,\Z)\cong 0\oplus\Hom(\pi_{3}(G),\Z)\cong\Z$；同理 $H^{2}(LG_{\C},\Z)\cong\Z$。末句把 (1.7) 的 $\omega$ 与其上同调类等同看待（本节已证 $\omega=\omega_{0}+d\beta$，$\omega_{0}$ 的类属于 $H^{2}(LG,\Z)$），并非笔误。\par\smallskip

\noindent \textbf{（原书 p.11）}\ 这里取 $\langle X,Y\rangle=\Tr(XY)$，于是 $\langle H,H\rangle=2$、$\langle E,F\rangle=1$。正是这组归一化使 $\{I_\mu\}$ 正规正交下的 $C=\sum_\mu I_\mu I_\mu$ 等于 $\frac12H^2+EF+FE$；下一段算出 Casimir 的特征值为 $2j(j+1)$，而 §1.2 末共形权公式里出现的 $j(j+1)$ 恰是它的一半，两处不要混。\par\smallskip

\noindent \textbf{（原书 p.13）}\ 原书 p.13 作 $u_m^*H=2mu_m$，右端漏了对偶基的星号，与同一句里「对偶基」的设定不符，疑为笔误；按配对 $\langle u_m^*,Hu_n\rangle=\langle u_m^*,2nu_n\rangle=2m\delta_{mn}$ 应为 $u_m^*H=2mu_m^*$。\par\smallskip

\noindent \textbf{（原书 p.14）}\ 这里 $V_\lambda$ 的维数是 $\lambda+1=2j+1$，同一段里 $\lambda$ 与 $j$ 两套记法并存（$\lambda=2j$）。命题 1.4 的条件 $\lambda_1+\lambda_2+\lambda_3\in2\Z$ 换回自旋就是 $j_1+j_2+j_3\in\Z$；此后凡见 $\lambda$ 都指最高权，勿与自旋 $j$ 或维数 $\lambda+1$ 混用。\par\smallskip

\noindent \textbf{（原书 p.14）}\ $A_+$ 的定义取 $n>0$（不含常数项），$A_-$ 取 $n<0$，故常数项 $\fg\otimes1$ 不属于 $\mathcal{N}_\pm$；$\mathcal{N}_0=\C H\oplus\C c$ 只补上零次的两个方向。正因如此，零次元 $X\otimes t^0$ 与 $H,c$ 是分开处理的，下面 (1.9)、(1.10) 才要按 $n=0$ 与 $n\ne0$ 分两种写法。\par\smallskip

\noindent \textbf{（原书 p.16）}\ (1.10) 中 $j$ 只跑 $j\ge1$，等于把 (1.9) 里 $j$ 与 $-j$ 成对的两项合并后只取一半，所以前因子从 $1$ 变成 $1/(k+2)$；若照 (1.9) 用 $j\in\Z$ 又保留 $1/(k+2)$，$L_0$ 会多算一倍。\par\smallskip

\noindent \textbf{（原书 p.16）}\ 这里把 Virasoro 交换关系写成 $[L_m,L_n]=(m-n)L_{m+n}+\frac{m^3-m}{12}\delta_{m+n,0}c$，其中 $c=3k/(k+2)$。$k\to\infty$ 时 $c\to3$，正是 $\fg=\fsl_2(\C)$ 的维数。$\delta_{m+n,0}$ 说明中心项只在 $n=-m$ 时出现，而 $m^3-m$ 在 $m=0,\pm1$ 处为零，故子代数 $\{L_{-1},L_0,L_1\}$ 上中心项不起作用。\par\smallskip

\noindent \textbf{（原书 p.17）}\ $\sigma$ 前的 $1/24\pi^2$ 是归一化因子，使 $\int_{SU(2)}\sigma=1$ 恰为 $H^3(SU(2),\Z)$ 的生成元；正因如此，WZW 项前面的 $\sqrt{-1}k/12\pi$ 才保证 $\exp(-S_\Sigma(f))$ 单值，$k$ 必须是整数。若把 $\sigma$ 的常数改成 $1/8\pi^2$（另一种常见约定），$k$ 的整性条件随之改变。\par\smallskip

\noindent \textbf{（原书 p.18）}\ 原书 p.18 把第二个扩张写成 $\tilde f':\Sigma\to G$，但同一句又称它是 $f$ 到 $B'$ 的扩张，定义域应为 $B'$；疑为笔误。证明的关键一步是「积分值必为整数」，即 $\sigma$ 的 de Rham 类属于整系数上同调，故 $\exp(-2\pi\sqrt{-1}\cdot\text{整数})=1$。\par\smallskip

\noindent \textbf{（原书 p.19）}\ 这一步的关键是 $I_2$ 的两个交叉项与 Wess--Zumino 项改写后的贡献合成一个全微分，故只剩 $I_1$。另请注意等式第三项只含 $f$ 的反全纯导数 $\bar\partial f$ 与 $g$ 的全纯导数 $\partial g$：当 $f$ 全纯或 $g$ 反正则时该项为零，这正是本节末命题 1.11 的出发点。\par\smallskip

\noindent \textbf{（原书 p.21）}\ 条件 (1) 只保证两条道路端点相同，能否拼成纤维上的同一元素要看条件 (2)；条件 (2) 右端正是把 $\exp(-S_{\CP^1}(h))$ 与和乐因子 $\Gamma_{D_\infty}$ 相比较，故说它是关于和乐的条件。$\mathcal{L}\cong L^{\otimes k}$ 说明此处的 $k$ 与 WZW 作用里的 level 是同一个整数。\par\smallskip

\noindent \textbf{（原书 p.22）}\ 配对定义的指数是正的 $\exp(S_{\CP^1}(f))$，与 $\exp(-S_\Sigma(f))$ 的负号相反：双对丛 $\mathcal{L}^{-1}$ 的作用正是把这个正号吃回去。也正因如此，(1.13) 的积里出现的是 $-\Gamma_D$ 而不是 $+\Gamma_D$。\par\smallskip

\noindent \textbf{（原书 p.24）}\ 证明只用了一条：$\Gamma_\Sigma(h,g)$ 的被积式含 $h^{-1}\bar\partial h$，$h$ 全纯时整个为零，于是 (1.14) 退化为 $\exp(-S_\Sigma(hg))$。反全纯的 $\bar h$ 恰好使 $\partial$ 一侧消失。故左作用只对全纯 $h$、右作用只对反全纯 $\bar h$ 协变，一般的 $h$ 只会剩下带 $\Gamma$ 的因子。\par\smallskip

\noindent \textbf{（原书 p.26）}\ 这里 $\fg(\bm p)$ 的元素取值为亚纯函数，故它在 $p_i$ 处只有有限阶极点；$\iota$ 把同一个 $f$ 在各点 $p_j$ 处分别展开、再直和起来。当 $f$ 无极点时，各 $\iota_j(f)$ 只有非负次项；若对每个 $j$ 都有 $\iota_j(f)=0$，则 $f$ 在各 $p_j$ 附近恒为 $0$，由解析延拓的唯一性得 $f\equiv0$，故 $\iota$ 为单射。\par\smallskip

\noindent \textbf{（原书 p.28）}\ 式 (1.15) 右端只对 $j\ne i$ 求和，是因为 $X\otimes t_i^{r}$ 只在 $p_i$ 有极点，在其余 $p_j$ 处是全纯的，其 Laurent 展开成为 $m\ge0$ 的 Taylor 级数；负幂项全部集中在左端的 $p_i$ 一栏。又注意 $r<0$ 使 $f$ 确在 $p_i$ 有极点，这是“移动极点”技巧的关键。\par\smallskip

\noindent \textbf{（原书 p.30）}\ (1.16) 中 $E$ 的幂 $m_j$ 全取非负、且总数恰为 $d_i=k-\lambda_i+1$，这正是零向量 $\chi_i$ 里 $E$ 的个数。证明时反复用 (1.15) 把 $t_1^{-1}$ 的极点逐个挪到 $p_2,\dots,p_n$，再比较 $\prod(z_j-z_1)^{-m_j}$ 的系数；这一“比较系数”之所以可行，是因为各 $z_j$ 与 $z_1$ 独立。\par\smallskip

\noindent \textbf{（原书 p.31）}\ 关键计算里 $H$ 的特征值是 $2(k+1)-(\lambda_1+\lambda_2+\lambda_3)+2(n_2+n_3)$：$+2(k+1)$ 来自 $V_{\lambda_1}$ 与 $d_1$ 的搭配，即 $d_1+\lambda_1=k+1$。由 $\lambda_1+\lambda_2+\lambda_3\in2\Z$，$2(k+1)-(\lambda_1+\lambda_2+\lambda_3)$ 与 $2(n_2+n_3)$ 同奇偶，故它为零当且仅当 $\lambda_1+\lambda_2+\lambda_3\ge2k+2$。这解释了“$\le 2k$ 时自动满足、$>2k$ 时给出独立条件”的分界。\par\smallskip

\noindent \textbf{（原书 p.32）}\ 两种维数记号要分清：三点情形写 $N_{\lambda_1\lambda_2\lambda_3}$（(1.18) 之前的 $\mathcal{H}(p_1,p_2,p_3;\cdot)$），而带上标的 $N^{\lambda_{n+1}}_{\lambda_1\cdots\lambda_n}$ 表示第 $n+1$ 个权来自对偶模 $V^{*}_{\lambda_{n+1}}$。$\mathfrak{sl}_2$ 的有限维不可约表示自对偶，故两种写法数值相同；换成一般单 Lie 代数时二者不可混用。\par\smallskip

\noindent \textbf{（原书 p.33）}\ 图中每条边上的箭头只标出该边的定向（相对顶点的指向），并不表示最高权的大小，真正约束每条边权的是三价顶点处的 level $k$ Clebsch--Gordan 条件 $\mu+\nu+\lambda\in2\Z$、$|\mu-\nu|\le\lambda\le\mu+\nu$、$\mu+\nu+\lambda\le2k$。中央与右侧两图给出同一空间的不同基，其间的变换就是下一章所说的辫子（braid）作用。\par\smallskip

\noindent \textbf{（原书 p.34）}\ 这个并集只用来指明纤维是 $\cH(p,\lambda)$，它本身还没有丛结构；真正的丛结构由下文把 $\cE_\lambda$ 嵌入平凡丛 $\Conf_n(\CP^1)\times\Hom_\C(\bigotimes_j\cH_{\lambda_j},\C)$ 来定义，不要把它当成简单的拓扑并集。\par\smallskip

\noindent \textbf{（原书 p.37）}\ 原书 p.37 此式左边求和对 $i$ 进行，而 $\Psi$ 中插入的是 $f_{D_i}\xi_i$，$L_{-1}$ 的上标却是 $(j)$，两处指标不一致；按下一行以 $g_{D_j}=\partial_{z_j}f_{D_j}+[L_{-1},f_{D_j}]$ 相消的推导，应以同一指标 $j$ 为准。此处照原书排印，存疑待考。\par\smallskip

\noindent \textbf{（原书 p.37）}\ 这里把联络写成 $\nabla=\rd-\omega$，$\omega$ 取值于 $\End$；曲率为 $\rd\omega-\omega\wedge\omega$。由于 $L_{-1}^{(i)}$ 与 $L_{-1}^{(j)}$ 作用在不同分量上，二者交换，$\omega\wedge\omega=0$ 是显然的，真正要算的是 $\rd\omega=0$，即 $L_{-1}^{(i)}$ 不显含 $z$。\par\smallskip

\noindent \textbf{（原书 p.39）}\ KZ 方程右端的系数 $1/(k+2)$ 来自 Sugawara 构造中 $L_{-1}$ 的归一化 (1.9)；若改动 $L_n$ 的定义或 $c$ 的归一化，这个系数随之改变。核对文献时先看其 $L_{-1}$ 的写法，再比对系数。\par\smallskip

\noindent \textbf{（原书 p.41）}\ 命题 1.19 把“新变量”写在右边：$\Psi_0(z)=\prod(cz_j+d)^{-2\Delta_j}\Psi_0(w)$，与 (1.29) 中 $\phi(v,f(z))=(f')^{-\Delta}\phi(v,z)$ 的方向一致；指数是 $2\Delta_{\lambda_j}$ 而不是 $\Delta_{\lambda_j}$，因为这里 $\Psi_0$ 有 $n$ 个协变分量、每个贡献一次权。\par\smallskip

\noindent \textbf{（原书 p.42）}\ 这里 $u$ 取自 $\cH_{\lambda_\infty}^{*}$ 而对偶空间的元素要做配对才落回 $\C$；把 $\phi(v,z)$ 读作“从 $\cH_{\lambda_0}$ 到 $\cH_{\lambda_\infty}$ 的算子”时，方向来自 $\cH^{*}$ 的自对偶与配对，不要误以为 $\Psi$ 的第三个因子 $\cH_{\lambda_0}$ 是作用的对象。\par\smallskip

\noindent \textbf{（原书 p.43）}\ (1.29) 中 $n$ 取遍整数，而 (1.24)、(1.25) 只对 $n\ge 0$ 陈述；两者说的是不同对象：前者是绕 $z$ 的围道积分给出的算子，后者是它在最高权表示顶端 $V_\lambda$ 上的分量。若把两处的 $n$ 混为一谈，会误以为负模作用为零。\par\smallskip

\noindent \textbf{（原书 p.45）}\ 图1.3 的实质是：先取满足 $|w|>|z|$ 的圆 $C_1$ 与满足 $|z|>|w|$ 的圆 $C_2$，两者之差等于绕 $z$ 的小圆 $C$；把 $L_n\phi$ 与 $\phi L_n$ 的围道积分相减，再用 (1.33) 的奇性项取留数，就得到 (1.29)。作者只写“取积分路”而省略了这一步留数计算。\par\smallskip

\noindent \textbf{（原书 p.46）}\ (1.37) 的系数 $3/2$ 是使 Schwarz 导数满足链式法则的唯一取法；若把它改掉，$S$ 就不再是共形不变量，(1.38) 的 $c/12$ 与 (1.34) 的 $c/2$ 也不能再与它配套。三处由同一个中心荷 $c$ 决定，改一处就必须同时改另外两处。\par\smallskip

\noindent \textbf{（原书 p.47）}\ (1.38) 与命题 1.21 的对应要看清自变量：$T(z)$ 换成 $T(w)$ 时乘的是 $(\partial w/\partial z)^{2}$，而 Schwarz 导数 $S(f,z)$ 的第三个变元仍记作 $z$。若把 $\partial w/\partial z$ 误写成 $\partial z/\partial w$，中心荷项的符号会跟着出错。\par\smallskip

\noindent \textbf{（原书 p.49）}\ (1.40) 里 $q^{L_0-c/24}$ 的 $c$ 是 Virasoro 中心荷，与 level $k$ 的关系为 $c=3k/(k+2)$；(1.43) 的 $\Delta_\lambda=\lambda(\lambda+2)/(4(k+2))$ 正是从 Sugawara 构造的 $L_0$ 在 $V_\lambda$ 上的特征值算出来的。把 $\Delta_\lambda$ 当成自由参数会与 (1.42) 的 $S$ 矩阵对不上。\par\smallskip

\noindent \textbf{（原书 p.50）}\ 命题 1.22 的右端以 $S_{0\alpha}$ 作分母，故 $\alpha$ 从 $0$ 取到 $k$；$\lambda,\mu,\nu$ 也同此范围。$N_{\lambda\mu\nu}$ 关于三个指标对称，这与 $R_k$ 的交换性一致，但 (1.22) 用的 $S$ 只对 $A_1^{(1)}$ 型给出 (1.42) 的显式公式，换其他型时 $S$ 要另行计算。\par\smallskip

\section*{Jones--Witten 理论}

\noindent \textbf{（原书 p.52）}\ 式 (2.1)、(2.2) 即 Artin 辫群的标准表示：相邻生成元之间是辫关系，不相邻生成元之间交换。第二式的条件 $|i-j|>1$ 表达的是“不相邻”，不是 $i\ne j$，两者常被混淆。\par\smallskip

\noindent \textbf{（原书 p.52）}\ 记号写作 $\rd\log(z_i-z_j)$，但 $\log$ 有分支，它没有单值的原函数；作为一次微分形式，它在 $\Conf_n(\C)$ 上整体定义且闭。又 $\omega_{ji}=\omega_{ij}$，故 (2.3) 三项是同一形式按 $i<j<k$ 循环取值的和。\par\smallskip

\noindent \textbf{（原书 p.53）}\ 参数 $\kappa$ 与共形场论的 level $k$ 由 $\kappa=k+h^\vee$ 相联系（$h^\vee$ 为对偶 Coxeter 数）；本书后面取 $\fg=\fsl_2(\C)$ 时 $h^\vee=2$，故 $\kappa=k+2$。与文献对比时这一步平移最易出错。\par\smallskip

\noindent \textbf{（原书 p.53）}\ 这里 $\{I_\mu\}$ 是 $\fg$ 关于 Cartan--Killing 型的标准正交基，$\Omega=\sum_\mu I_\mu\otimes I_\mu$ 与基的选取无关，是 Casimir 元在张量积上的作用；$\Omega_{ij}$ 只在前述两个因子上非平凡，其余因子取恒等映射。\par\smallskip

\noindent \textbf{（原书 p.54）}\ 平坦性 = 曲率为零，即 $\rd\omega+\omega\wedge\omega=0$；此处 $\rd\omega=0$ 与 (2.9) 各贡献一项。注意 $\omega$ 是矩阵值一次形式，$\rd\omega$ 与 $\omega\wedge\omega$ 的次序不可随意交换。\par\smallskip

\noindent \textbf{（原书 p.54）}\ 解按列排成 $m\times m$ 矩阵 $(\Phi_1,\dots,\Phi_m)$，单值矩阵右乘；故 (2.16) 中 $F$ 的左乘、右乘位置要与 $\Phi_1,\Phi_2$ 的排法对照，不能只记式子。\par\smallskip

\noindent \textbf{（原书 p.55）}\ $\Delta_\lambda$ 是最高权 $\lambda$ 对应的共形权（Virasoro 最高权）。式 (2.12) 的第一个值来自 $\Omega_{12}$ 的作用，第二个值来自 $\Omega_{12}+\Omega_{13}+\Omega_{23}$ 的作用；两者都是 Sugawara 构造给出的 Virasoro 作用，故只依赖共形权之差。\par\smallskip

\noindent \textbf{（原书 p.57）}\ 原书 p.57 式 (2.19) 作 $G_2$，但 $G_2$ 全书未定义；按上下文（$G_0$ 在 $x=0$ 附近、$G_1$ 在 $x=1$ 附近，解析延拓给出二者的关系）此处应为 $G_0$。正文照录原书。\par\smallskip

\noindent \textbf{（原书 p.58）}\ 树形图基的关键在于：对每个 $k$，$\sum_{1\le i<j\le k}\Omega_{ij}$ 在该基下同时对角化。因此 (2.20)、(2.21) 中 $u_k$、$v_k$ 的指数可以逐项写出而不出现混合项，这也是后续把所有连接矩阵归结到 $n=3$ 的依据。\par\smallskip

\noindent \textbf{（原书 p.59）}\ 命题 2.2 就是结合子满足的五边形公理：两种“先合并哪一对”的路径给出同一连接矩阵。把图 2.4 的图换成它们的对偶图即得标准的 pentagon，这也是 Drinfel'd associator 定出辫群表示的原因。\par\smallskip

\noindent \textbf{（原书 p.61）}\ $P_{12}$ 是张量积因子置换算子，$\Omega_{12}$ 是 Casimir 作用；二者交换，故指数量 $P_{12}\exp(\cdots\Omega_{12})$ 可按 $\Omega_{12}$ 的特征分解逐项计算。这里的 $\sqrt{-1}$ 与 $\kappa$ 的配合决定了固有值是否为实数。\par\smallskip

\noindent \textbf{（原书 p.61）}\ 引入 $\eta$ 是为把 $\rho(\sigma_i)$ 的特征值写成 $q$ 的幂次的对称形状；由 (2.23) 直接看出 $\tilde\rho(\sigma_i)=g_i$ 已满足 Iwahori--Hecke 代数的二次关系 $(T_i-q^{1/2})(T_i+q^{-1/2})=0$。\par\smallskip

\noindent \textbf{（原书 p.62）}\ 定理 2.3 的证明在此只用“归纳上面的考察”一句带过。实际上需要验证 $\tilde\rho$ 满足 (2.1)、(2.2) 与 (2.23) 三组关系；其中辫关系由 KZ 联络在 $\Conf_n(\C)$ 上的平坦性保证（同伦不变），二次关系由 (2.23) 直接给出。\par\smallskip

\noindent \textbf{（原书 p.64）}\ 图 2.7 中正负交叉的约定依赖两条分支的定向：把其中一条分支的定向反转，同一交叉便由正变负。$S^3$ 中链环的环绕数必须先固定各分支定向才有定义。\par\smallskip

\noindent \textbf{（原书 p.66）}\ 图 2.10 的三种变形分别对应辫关系 $\sigma_i\sigma_{i+1}\sigma_i=\sigma_{i+1}\sigma_i\sigma_{i+1}$、互不相邻生成元的交换，以及极大点与极小点沿曲线的移动；引理的实质就是 KZ 连接的平坦性。\par\smallskip

\noindent \textbf{（原书 p.67）}\ 归一化所用的是极大点个数 $\mu_j$，而不是「极大点个数减极小点个数」。任何链环分支的投影图上极大点与极小点成对出现、个数相等，故两种计数给出的幂次相同；若把图 2.11 的约定改成以极小点归一，$d(\lambda)$ 的幂次会整体变号。\par\smallskip

\noindent \textbf{（原书 p.67）}\ 图 2.12 的三幅图给出 $L_+$、$L_-$、$L_0$ 的局部替换：圆内分别是正交叉、负交叉与两条不相交的弧，圆外部分完全固定。拆接关系式比较的正是这三者在同一位置的不同。\par\smallskip

\noindent \textbf{（原书 p.68）}\ 此处的 $q$ 与 §2.1 的 $q$ 差一个 $m$ 次根：$q^{1/m}=\exp(2\pi\sqrt{-1}/(m(k+2)))$。拆接关系式含 $q^{\pm1/4}$、$q^{\pm1/2}$，故必须先固定 $q^{1/m}$ 的分支；$q$ 为单位根这一事实使 $J_L$ 良定义，也使它化为 Jones 多项式的特殊值。\par\smallskip

\noindent \textbf{（原书 p.69）}\ $2$ 维共形块空间中的三个线性映射必线性相关，这就是拆接关系式的来源；系数 $q^{\pm1/4}$、$q^{\pm1/2}$ 由命题 1.15 的融合规则（即 $F$ 矩阵）唯一确定，并非任意选取。\par\smallskip

\noindent \textbf{（原书 p.70）}\ 图 2.14 的 $L_0$ 是两个互不相交的平凡纽结，其不变量并不等于 $d(1)$ 本身：由命题 2.10，两分支权皆为 $1$ 时 $J_{L_0}=d(1)^2$，而 $P_L$ 还含因子 $d(1)^{-1}$，故 $P_{L_0}=d(1)$。原书略去了这一步。\par\smallskip

\noindent \textbf{（原书 p.71）}\ 原书此式右端写作 $J(K;\nu)$，而 $K$ 在前文并未定义（cabling 的两个分支记作 $K_0,K_1$）。按上下文与引理 2.9 的用法，$K$ 应理解为本来的结 $K_0$；此处照原书抄录，未擅改。\par\smallskip

\noindent \textbf{（原书 p.72）}\ 这里 $d(\lambda)$ 有两个来源：一处由 $J(K;\lambda)$ 定义，一处由 $S_{0\lambda}/S_{00}$ 给出。证明的思路是先验证后者满足式 (2.26)、再用「式 (2.26) 加 $d(1)$ 的值唯一确定 $d(\lambda)$」把两者接上；Verlinde 等式正是用来验证式 (2.26) 的。\par\smallskip

\noindent \textbf{（原书 p.72）}\ $q$ 为单位根时 $q^{-1}=\bar q$，所以 $J_L(q^{-1})=\overline{J_L(q)}$；若 $q$ 不是单位根，这一步不成立。这正是本章只在 $q$ 为单位根（即 $k$ 为正整数）时讨论不变量的原因。\par\smallskip

\noindent \textbf{（原书 p.73）}\ 图 2.16 中各值均为 Laurent 多项式：trefoil 为 $t+t^3-t^4$，把 $t$ 换成 $t^{-1}$ 得镜像的值 $t^{-1}+t^{-3}-t^{-4}$，正是本段所说的 $V_{\bar L}(t)=V_L(t^{-1})$。虚线圆标出每次作拆接的位置。\par\smallskip

\noindent \textbf{（原书 p.74）}\ $\widehat{T_1\circ T_2}$ 的帽子覆盖整个合成 $T_1\circ T_2$，表示把合成后的 tangle 整体闭合；$\hat T_1,\hat T_2$ 则是分别闭合。两者是不同的链环，右端把两个闭合结果的不变量相乘，因子 $S_{00}/S_{0\lambda}$ 来自闭合处权 $\lambda$ 的收缩。\par\smallskip

\noindent \textbf{（原书 p.75）}\ 对 $T_+$ 只看「$1$ 的像」，是因为 $T_+$ 的 tangle 算子是 $\C=V(t_0)\to V(t)$ 的映射；$T_-$ 上下翻转后其算子是 $V(t)\to\C$，故 $J(T_-)$ 落在对偶空间 $V(t)^\ast$ 中，与 $J(T_+)$ 配对即得标量 $J(L)$。\par\smallskip

\noindent \textbf{（原书 p.76）}\ $E(L)=S^3\setminus N(L)$ 是链环的外部，其边界是 $m$ 个环面；$f$ 把实心环面的经线送到 framing 曲线 $\gamma_j$。因此这里的「framing」并不是额外的装饰，它等价于粘合映射 $f$ 的选取。\par\smallskip

\noindent \textbf{（原书 p.76）}\ 平凡纽结作 framing $n$ 的 Dehn 手术得透镜空间 $L(n,1)$：$n=\pm1$ 时为 $S^3$，$n=0$ 时为 $S^2\times S^1$。可见同一三维流形可由不同的链环得到，这正是下一节引入 Kirby 操作的缘由。\par\smallskip

\noindent \textbf{（原书 p.76）}\ $L\subset\partial B^4$ 的 framing 决定 $2$-柄的粘法；$W_L$ 的构造说明 Dehn 手术是四维柄分解的一部分，$M_L=\partial W_L$ 把三维流形的同胚问题化为链环的 Kirby 等价问题。\par\smallskip

\noindent \textbf{（原书 p.77）}\ 对 $i\ne j$，$L_i\cdot L_j$ 是两条不同分支的环绕数；对角元 $L_i\cdot L_i$ 则是确定 $L_i$ 的 framing 的整数，须理解为 $L_i$ 与其平行副本的配对，并不是 $L_i$ 与自身的环绕数。\par\smallskip

\noindent \textbf{（原书 p.78）}\ 原书此处的指数是 $-3$，不是 $-1$：$C=\exp(2\pi\sqrt{-1}c/24)^{-3}$。后文定理 2.16 中的 $C^{\sigma(L)}$ 用的就是这个量，转抄时不要漏掉指数。\par\smallskip

\noindent \textbf{（原书 p.79）}\ 引理 2.14 中 $\lambda$、$\nu$ 是固定的外指标，求和只对中间的 $\mu$ 进行；指数里三个 $\Delta$ 的下标分别为 $\lambda,\mu,\nu$，不要漏掉 $\Delta_\mu$。\par\smallskip

\noindent \textbf{（原书 p.81）}\ 证明中 $\sigma(L)$ 与 $\sigma(L')$ 之差不是零，而是被 $C^{\sigma(L)}$、$C^{\sigma(L')}$ 两个因子吸收；如果误以为 Kirby 操作保持符号差，就会看不出这一步为何成立。\par\smallskip

\noindent \textbf{（原书 p.81）}\ 这两个例子给出的是归一化基准：$S^3$ 的值取 $S_{00}$ 而非 $1$，$S^1\times S^2$ 的值恰好为 $1$。后文命题 2.17 中的因子 $1/S_{00}$ 正是为了抵消 $S^3$ 这一归一化。\par\smallskip

\noindent \textbf{（原书 p.82）}\ 命题 2.18 中右端的横线是复共轭，不是闭包也不是取逆；它与命题 2.11（镜像与复共轭的对应）是同一事实的两种说法。\par\smallskip

\noindent \textbf{（原书 p.83）}\ (A1)、(A2) 决定了“边界反定向”与“不交并”的处理，由此才有 $Z(Y)\in\Hom_\C(Z_{\Sigma_1},Z_{\Sigma_2})$：若把 $\partial Y=(-\Sigma_1)\sqcup\Sigma_2$ 误写成 $\Sigma_1\sqcup\Sigma_2$，映射的方向就会反过来。\par\smallskip

\noindent \textbf{（原书 p.84）}\ 取下 $f=\id$ 时 $\rho(f)$ 是恒等映射，取迹即得 $\dim Z_\Sigma$；这正是二维拓扑场论里“配分函数等于态空间维数”的来源，也说明 $Z(\Sigma\times S^1)$ 必须按无边界情形解释。\par\smallskip

\noindent \textbf{（原书 p.85）}\ $V_\Sigma$ 是所有允许权对应的 $V_{\mu_1\mu_1^*\cdots\mu_g\mu_g^*}$ 的直和，不是其中某一项；下标中 $\mu_i$ 与 $\mu_i^*$ 成对出现，正对应图 $\Gamma$ 的两端各一个边界点。\par\smallskip

\noindent \textbf{（原书 p.88）}\ 把 $h$ 换一种 Lickorish 生成元的分解方式时，$\rho(h)$ 只会相差一个标量，这个标量正是 $\xi$；因此 $\rho$ 是射影表示而不是线性表示，不能直接当成群同态使用。\par\smallskip

\noindent \textbf{（原书 p.89）}\ 式 (2.28) 中的 $S$ 既是 $SL_2(\Z)$ 的生成元，又是 $S$ 矩阵的来源，两者记号相同但含义不同；指标顺序为 $S v_\lambda=\sum_\mu S_{\lambda\mu}v_\mu$，其中行指标是 $\lambda$，列指标是 $\mu$。\par\smallskip

\noindent \textbf{（原书 p.90）}\ 命题 2.20 的因子 $S_{00}^{-g+1}$ 正是为了让上式中的 $S_{00}$ 被吸收；若漏掉这个幂次，稳定等价下两边会差一个 $S_{00}$ 的倍数。\par\smallskip

\noindent \textbf{（原书 p.91）}\ 添加 1-柄的算子 $\iota$ 前面的 $1/\sqrt{S_{00}}$ 与球面生成元 $\beta$ 的归一化是同一约定；漏掉这个因子会使公理 (A3) 的偏差不再是 $C$ 的幂次。\par\smallskip

\noindent \textbf{（原书 p.92）}\ (2.30) 只在 $M=S^3$ 时是构造中选定的归一化条件，并非对一般 $M$ 成立；一般 $M$ 的值须回到定理 2.16 的 Dehn 手术公式计算。\par\smallskip

\noindent \textbf{（原书 p.93）}\ 此处的 $C$ 是前文 $Z_k(M)$ 归一化中出现的同一个符号，并非复数域 $\C$；“以 $C$ 的幂为模”意为结果只在相差 $C^{n}$ 的意义下相等。\par\smallskip

\noindent \textbf{（原书 p.93）}\ $S_{00}J(L;\lambda)$ 是 $S$ 矩阵的 $(0,0)$ 元与 Jones 不变量之积；framing $0$ 指取零标架，改换标架会改变该值，故不可省略。\par\smallskip

\noindent \textbf{（原书 p.95）}\ 式 (2.34) 中的 $\rd$ 是无限维空间 $\cA_M$ 上的微分（变分），不是 $M$ 上的外微分；系数 $1/4\pi^{2}$ 与 (2.33) 的 $1/8\pi^{2}$ 差因子 2，来自 $\delta(A\wedge\rd A)$ 两项的合并，可与命题 2.21 的展开式对照核对。\par\smallskip

\noindent \textbf{（原书 p.95）}\ (2.36) 说明规范变换只把 $CS$ 改变一个整数值，这正是 $e^{2\pi\sqrt{-1}k\,CS}$ 在整数 $k$ 下规范不变的根源，也说明 $Z_k(M)$ 为何只对整数 $k$ 有定义。\par\smallskip

\noindent \textbf{（原书 p.96）}\ 原书定义 $L_{\Sigma,a}$ 时下标用一般元素 $a$，其后一律写作 $\alpha$（$\alpha$ 是 $A$ 到 $\Sigma$ 的限制）；两者是同一类对象，抄录时保留原样，勿据此后退修改公式。\par\smallskip

\noindent \textbf{（原书 p.97）}\ 把 $\exp(2\pi\sqrt{-1}CS)$ 看成线丛的截面而非函数，是因为它依赖 $a$ 到 $M$ 的扩张 $A$；沿 $\cA_\Sigma$ 的道路走一圈（闭路）时按 $c(a,g)$ 作和乐平移，这正是“截面”而非“函数”的含义。\par\smallskip

\noindent \textbf{（原书 p.98）}\ 原书 p.98 此处印作 $c(g,\alpha)$，与 (2.38) 中 $c(a,g)$ 的自变量次序相反；按定义（$a$ 取在 $\cA_\Sigma$，$g$ 取在规范群）应为 $c(\alpha,g)$。疑为笔误，正文照录原文。\par\smallskip

\noindent \textbf{（原书 p.98）}\ Kähler 极化指在 $\cM_\Sigma$ 上取定与辛形式相容的复结构；不同的极化给出同一空间的不同模型（正则是“全纯截面”相对哪个复结构而言），作者未讨论其依赖，引用文献 [19] 处理。\par\smallskip

\noindent \textbf{（原书 p.98）}\ (2.41) 中的路径积分是形式记号，$\mathcal{D}A$ 没有测度意义；本节所有含 $\int\mathcal{D}A$ 的等式都只在微扰展开（第 3 章）的意义下理解。\par\smallskip

\noindent \textbf{（原书 p.99）}\ 原书《要約》只列 2.1--2.4 四条，没有 §2.5，这是原书面貌，不是漏译；§2.5 的内容（联络空间的辛几何与复线丛）可视为 2.3、2.4 的技术准备。\par\smallskip

\noindent \textbf{（原书 p.99）}\ 2.1 条的“特殊情形”指 $G=SU(m)$、$V$ 取 $m$ 维自然表示：此时 KZ 表示经由 Iwahori--Hecke 代数分解；一般情形的单值群并不落在 Hecke 代数像中。\par\smallskip

\noindent \textbf{（原书 p.99）}\ 2.2 条的“1 的单位根”即 $q=e^{2\pi\sqrt{-1}/k}$ 型单位根；此时 $U_q(\fsl_2)$ 的表示论给出 Jones 多项式的特殊值，这是量子群与共形场论相通的入口。\par\smallskip

\section*{Chern--Simons 微扰理论}

\noindent \textbf{（原书 p.110）}\ 本节把三个层次分得很清楚：$V_m$（纽结不变量）$\to$ $\cA_m$（弦图模 4 项关系与孤立弦关系）$\to$ 权系统（$\cA_m^{*}$ 的元素）。定理 3.3 说 Kontsevich 积分给出 $V_m \to \cA_m^{*}$ 的逆映射，所以「所有 Vassiliev 不变量」与「所有权系统」一一对应，这是 §3.2 构造 $Z_k(M)$ 渐近展开的出发点。\par\smallskip

\noindent \textbf{（原书 p.110）}\ (3.1) 的求和约定容易读错：$\varepsilon_j$ 取 $\pm1$ 共有 $2^{n}$ 项，$\varepsilon_1\cdots\varepsilon_n$ 是符号因子；$n>m$ 时恒为零正是「$m$ 阶不变量对 $m+1$ 个二重点的符号和为零」，与 (3.2) 的 4 项关系式是同一件事的两种写法。\par\smallskip

\noindent \textbf{（原书 p.110）}\ (3.17) 的指数是 $-\mu(K)+1$（$\mu(K)$ 为极大点个数），不是 $-\mu(K)-1$。这是因为 $(2\pi\sqrt{-1})^{-m}$ 的归一化与 $K_0$ 的极大点计数一起吸收掉了 framing 与临界点对消的差异；页图 p.108 已逐字复核。\par\smallskip

\noindent \textbf{（原书 p.110）}\ (3.15) 的 $D_P$ 只在「$D_P$ 等于图 3.2 右侧的图」时取 1，它是把弦图映为 $\cA_m$ 中元素后的系数；$\varepsilon(P)$ 数的是虚线的端点中「$K$ 的定向朝下」的个数，方向颠倒会整体变号，核对时以页图 p.107 的图为准。\par\smallskip

\noindent \textbf{（原书 p.110）}\ $Z(K_0)$ 的公式里 $m$ 同时出现在分子（$m(q^{1/2}-q^{-1/2})$）与分母的指数（$q^{m/2}-q^{-m/2}$），其中分子的 $m$ 来自 $m$ 维表示 $\mathfrak{sl}_m$ 的维数，不要与弦的条数 $m$ 混淆（本例中两者恰好同号）。\par\smallskip

\noindent \textbf{（原书 p.111）}\ 这里的归一化系数 $1/8\pi^{2}$ 与 $2/3$ 是配套的：对闭的 $M$、$G=SU(2)$，$CS(g^{*}A)-CS(A)$ 恰取整数值 $n$，$n$ 是 $g\colon M\to SU(2)$ 的映射度。对一般的单 Lie 群，系数须按 Killing 型的长根归一化相应调整，才能保持 $n\in\mathbb{Z}$。\par\smallskip

\noindent \textbf{（原书 p.112）}\ 这里的 $\fg_{\alpha}$ 是“局部系统”，即以 $\pi_{1}(M)$ 的表示 $\rho_{\alpha}$ 为纤维、在 $M$ 上扭成的平坦向量丛的截面层；$H^{j}(M,\fg_{\alpha})$ 是局部系统数上同调（在 de Rham 意义下即用 $d_{\alpha}$ 定义的上同调）。$H^{0}=0$ 等价于 $\rho_{\alpha}$ 不可约，$H^{1}=0$ 等价于 $\alpha$ 在模去规范等价后的平坦联络模空间中是孤立点，这两条保证后面的算子行列式都有限。\par\smallskip

\noindent \textbf{（原书 p.113）}\ (3.22) 的相位 $e^{\frac{\pi i}{4}\operatorname{sgn} H(f,\alpha)}$ 就是平稳相位法的 Morse 指标相位：$\operatorname{sgn} H$ 等于 Hessian 的负特征值个数减去正特征值个数，故每个临界点贡献 $e^{\pm i\pi/4}$。因子 $k^{-n/2}$ 说明主项相当于“单圈”贡献；一元情形 $f''(x_{0})>0$ 时相位为 $e^{i\pi/4}$，与 (3.20) 中 $\lambda>0$ 的情形一致。\par\smallskip

\noindent \textbf{（原书 p.114）}\ 在三维流形上，$d^{*}_{\alpha}$ 与 $d_{\alpha}$ 通过 Hodge 星号相联系：在 $1$ 次形式上 $d^{*}_{\alpha}=-\ast d_{\alpha}\ast$（符号随次数与度规的号差而定）。正因如此，$Q=-*d_{\alpha}$ 才能写成 (3.23) 中的右上块，而 $P$ 是自伴算子，$P^{2}$ 才能按 Hodge 分解为两块 Laplace 算子的直和。$\varepsilon$ 的符号正是为了让 $P$ 自伴。\par\smallskip

\noindent \textbf{（原书 p.116）}\ $\Delta^{0}_{\alpha}=d^{*}_{\alpha}d_{\alpha}$ 是因为 $\Omega^{0}_{\alpha}$ 上没有 $(-1)$ 次形式，$d_{\alpha}d^{*}_{\alpha}$ 这一项在 $0$ 次上恒为 $0$；指数 $3/2$ 与 $1/2$ 则来自 $\ast:\Omega^{0}\to\Omega^{3}$、$\ast:\Omega^{1}\to\Omega^{2}$ 的配对------在三维流形上 Hodge 星号把一个 $0$ 次形式变成 $3$ 次形式，故 $\Delta^{0}$ 与 $\Delta^{3}$ 的谱相同，这正是 $(\det\Delta^{0})^{3/2}$ 中 $3$ 的来源。\par\smallskip

\noindent \textbf{（原书 p.116）}\ 指数中的 $k+2$ 相当于把 level 平移 $k+h^{\vee}$，$h^{\vee}=2$ 正是 $SU(2)$ 的 dual Coxeter 数；这一平移与 2-framing 的选取有关。$I_{\alpha}$ 只模 $8$ 确定，而 $e^{-2\pi i(I_{\alpha}+8)/4}=e^{-2\pi iI_{\alpha}/4}$，故相位无歧义。\par\smallskip

\noindent \textbf{（原书 p.125）}\ (3.25) 中 $k/4\pi$ 的归一化与第 2 章 (2.41) 一致；$\int_{K_i}A$ 两项是 Wilson 线在 $U(1)$ 情形的线性项，正是它们把链环信息带进配分函数。若把 $k$ 换成 $1/\hbar$ 并与 (3.27) 对照，可以看出「大 $k$ 渐近」与「半经典极限」是同一件事。\par\smallskip

\noindent \textbf{（原书 p.125）}\ (3.26) 的 $w(x)$ 是 $\R^{3}$ 中球面 $S^{2}$ 的体积元素在 $\R^{3}\setminus\{0\}$ 上的拉回，$\int_{S^{2}}w=1$ 的归一化决定了 Green 形式 $L(x,y)=w(x-y)$ 与环绕数积分 $\int\!\!\int w(x-y)$ 的符号约定；符号一旦改，整个 $I_{\Gamma}$ 的符号随之改变。\par\smallskip

\noindent \textbf{（原书 p.125）}\ (3.28) 中 $m=2$ 的两项分别对应两种把 4 个指标两两配对的缩并方式（$\lambda^{ii'}\lambda^{jj'}\lambda^{kk'}$ 与 $\lambda^{ij}\lambda^{kk'}\lambda^{i'j'}$），在图上就是 2 圈 Feynman 图的两种连边方式。\par\smallskip

\noindent \textbf{（原书 p.125）}\ 定理 3.5 中的 $(1/4)CS_{\mathrm{grav}}(g,s)$ 是「引力 Chern--Simons 项」，依赖切丛的 framing $s$ 的选取；它补偿了 $I_{\Gamma_1}+I_{\Gamma_2}$ 在改变 framing 时的变化，因此最终结果才可能不依赖度量的选取。作者在正文中略去了这一项的构造，只说明它出现。\par\smallskip

\noindent \textbf{（原书 p.125）}\ (3.32) 与 (3.33) 都含有 $\sum_a\gamma_{acd}\gamma_{aef}$ 型的结构常数缩并；用 Jacobi 等式 $\gamma_{ace}\gamma_{bfd}$ 型恒等式与反对称性把它们消去，是证明「$\nabla I_{\Gamma}$ 只有边界项」的关键一步，也是全书唯一一处真正使用 Lie 代数 Jacobi 等式的计算。\par\smallskip

\section*{今后的方向与课题}

\noindent \textbf{（原书 p.127）}\ 所谓 Abel 化，是把原本要靠非 Abel 对象（共形块、射影联络的和乐）描述的量，搬到含阿贝尔簇（Prym 簇）的 theta 函数语言中去；二重分支覆盖正是把模空间问题化为 Prym 簇上 theta 函数问题的桥梁。\par\smallskip

\noindent \textbf{（原书 p.128）}\ 记号 $\tau$ 在本书中指向多种互不相同的对象：p.128 的 $\tau_r(M)$、$\tau(M)$ 是大槻型不变量，p.131 的 $\tau$ 是 $A(\Sigma)\to C^{\infty}(M_{\Sigma})$ 的映射，p.133 的 $\tau$ 函数则属于 KdV 方程；须按上下文区分，勿当成同一记号。\par\smallskip

\noindent \textbf{（原书 p.130）}\ 式中的 $1/24$ 是 Chern--Simons 微扰计算给出的常数项（framing 修正的一部分），与 §1.1 式 (1.8) 的 $1/24\pi^{2}$、式 (1.7) 的 $1/4\pi^{2}$ 没有直接关系；不可因数字相同而把它们对应起来。\par\smallskip

\noindent \textbf{（原书 p.131）}\ 4 项关系式即弦图的 4T 关系；按它作商之后 $A(\Sigma)$ 上才可能有 Poisson 结构，这与第 3 章用弦图定义 Vassiliev 不变量时所用的关系式是同一套。\par\smallskip

\noindent \textbf{（原书 p.132）}\ 展开式中 $N^{2-2g}$ 的指数就是亏格 $g$ 闭曲面的 Euler 示性数：每增加一个把手，$N$ 的幂次降 2。这正是把矩阵积分读成曲面计数（大 $N$ 展开）的机制。\par\smallskip

\noindent \textbf{（原书 p.132）}\ 这里的「二次微分」是 quadratic differential，即 $dz^{2}$ 型的张量，与 2 次微分形式（2-form）不是同一对象；Mumford--Harer 的胞腔分解正是用它（Strebel 微分）的水平轨线作出的。\par\smallskip
